\section*{A060: the optimal constant-branch selector tax}

All cubes are uniform, Fourier coefficients are normalized, and logarithms
are base two.  Let (D:\{-1,1\}^n\to\{-1,1\}) be arbitrary and introduce
one fresh selector (x).  Define
\[
 K(+1,z)=1,\qquad K(-1,z)=D(z).
\]
Put
\[
 p=\Pr[D=-1]=\frac{1-\mathbb ED}{2}.
\]

\subsection*{Exact bias reduction}

\paragraph{Theorem (PROVED).}
For every (n\geq0) and every Boolean (D),
\[
 I_x(K)=p,\qquad I(K)=\frac12I(D)+p,
\]
and
\[
 H(K)=\frac12H(D)+p\,c(p),
\]
where, for (0<p<1),
\begin{align*}
 c(p)=\frac1p\bigg[&
 2(1-p)^2\log_2\frac1{1-p}
 +2p^2\log_2\frac1p+4p(1-p)\\
 &-(1-2p)^2\log_2\frac1{|1-2p|}
 \bigg],
\end{align*}
with the last term interpreted as zero at (p=1/2), and with the endpoint
values understood by continuity.  In particular, the entire selector tax
(H(K)-H(D)/2) depends only on the bias (p), not on any other Fourier
geometry of (D).

\paragraph{Proof.}
Write
\[
 K=u+xv,\qquad u=\frac{1+D}{2},\quad v=\frac{1-D}{2}.
\]
If (m=\widehat D(\varnothing)=1-2p), then the two constant-frequency
squared masses of (K) are
\[
 \widehat u(\varnothing)^2=(1-p)^2,
 \qquad
 \widehat v(\varnothing)^2=p^2.
\]
For every nonempty (S\subseteq[n]), the coefficient
(\widehat D(S)) produces the two squared masses
\[
 \widehat u(S)^2=\widehat v(S)^2=rac14\widehat D(S)^2.
\]
The entropy contribution of this pair is
\[
 \frac12\widehat D(S)^2
 \log_2\frac1{\widehat D(S)^2}
 +\widehat D(S)^2.
\]
Summing over nonempty (S), using
(\sum_{S\ne\varnothing}\widehat D(S)^2=1-m^2=4p(1-p)), and subtracting
half of the constant-frequency entropy of (D) gives exactly the displayed
formula for (p c(p)).

Characters containing (x) carry precisely the Fourier mass of (v), so
\[
 I_x(K)=\sum_S\widehat v(S)^2=\mathbb E v^2=p.
\]
Every nonempty data character of (D) appears twice with one quarter of its
old squared mass, hence contributes half of its old degree; the (x)-degree
adds (p).  Thus (I(K)=I(D)/2+p). \qed

\subsection*{Sharp universal coefficient}

Define
\[
 C_*=\max_{0\leq p\leq1}c(p).
\]

\paragraph{Theorem (PROVED).}
The function (c) has a unique global maximizer (p_*\in(0,1/2)).  It is
the unique zero in that interval of
\begin{align*}
 F(p)={}&2(1-p^2)\ln(1-p)-2p^2\ln p-4p^2\ln2\\
 &-(1-4p^2)\ln(1-2p).
\end{align*}
Moreover,
\[
 0.22944127<p_*<0.22944128,
 \qquad
 4.872321<C_*<4.872322.
\]
Consequently (C_*) is the smallest constant for which
\[
 H(K)\leq\frac12H(D)+C_* I_x(K)
\]
holds for every finite-dimensional Boolean (D).

\paragraph{Proof.}
Let (A(p)=(\ln2)\,p c(p)).  For (p\ne1/2),
\[
 A(p)=-2(1-p)^2\ln(1-p)-2p^2\ln p
 +4(\ln2)p(1-p)+(1-2p)^2\ln|1-2p|.
\]
Since (c'(p)=F(p)/(p^2\ln2)), it suffices to determine the sign of
(F).  Direct differentiation gives
\[
 F'(p)=pA''(p)
 =4p\ln\frac{(1-2p)^2}{4p(1-p)}.
\]
The logarithm changes sign only at
\[
 \alpha=\frac{1-1/\sqrt2}{2},\qquad
 \beta=\frac{1+1/\sqrt2}{2}.
\]
Also (F(0+)=0), (F(1/2)=-2\ln2), and (F(1-)=-4\ln2).
Thus (F) first increases on ((0,\alpha)), then strictly decreases to
(p=\beta), and finally increases while remaining negative up to (p=1).
It therefore has exactly one zero (p_*\in(\alpha,1/2)); (c) increases
before (p_*) and decreases afterward.

The attached rational-interval certificate proves
(F(0.22944127)>0), (F(0.22944128)<0), and evaluates (c) on that
root box.  A Boolean witness at (p=235/1024) proves
(C_*>4.872321), while interval evaluation on the root box proves
(C_*<4.872322).

Every Boolean density is dyadic, but dyadic numbers are dense.  For every
dyadic (p=k/2^n), a Boolean (D) of that density exists, and the exact bias
identity gives tax ratio (c(p)).  Dyadic densities approaching (p_*)
therefore approach (C_*).  No smaller universal coefficient can work.
\qed

\subsection*{Concrete near-extremal Boolean gate}

Index the ten-cube lexicographically by (0,\ldots,1023), and let
(D_{10}=-1) on indices (0,\ldots,234) and (+1) elsewhere.  Then
\[
 p=\frac{235}{1024},\qquad
 \mathbb ED_{10}=\frac{277}{512},\qquad
 I(D_{10})=\frac{277}{256}.
\]
The exact absolute integer-Walsh histogram of (D_{10}) is
\[
 \{2^{(768)},6^{(128)},10^{(64)},22^{(32)},42^{(28)},
 470^{(3)},554^{(1)}\},
\]
where coefficients are divided by (1024).  For its eleven-bit splice
(K_{11}),
\[
 I_x(K_{11})=\frac{235}{1024},\qquad
 I(K_{11})=\frac{789}{1024},\qquad
 \mathbb EK_{11}=\frac{789}{1024},
\]
and the absolute integer-Walsh histogram is
\[
 \{2^{(1536)},6^{(256)},10^{(128)},22^{(64)},42^{(56)},
 470^{(7)},1578^{(1)}\},
\]
with coefficients divided by (2048).  Rigorous interval evaluation yields
\[
 \frac{H(K_{11})-H(D_{10})/2}{I_x(K_{11})}
 =4.8723210430423538411\ldots>4.872321,
\]
whereas
\[
 \frac{H(K_{11})}{I(K_{11})}
 =3.2266979050419697025\ldots.
\]
Disjoint tensor powers preserve the latter ratio, not the local tax
coefficient.

\subsection*{FEI closure consequence and scope}

If a class containing (D) obeys (H(D)\leq C I(D)), then
\[
 H(K)\leq\frac C2 I(D)+C_*p
 \leq\max\{C,C_*\}\left(\frac12I(D)+p\right)
 =\max\{C,C_*\}I(K).
\]
Thus constant-branch fresh-selector extension preserves FEI with the sharp
scalar overhead (C_*).  This is a restricted closure theorem, not full FEI.
It keeps half of the child entropy and half of the child influence in the
ledger; it does not claim that selector influence alone pays the whole gate.

Because the bias reduction is exact for every child (D), it includes rare
defects and addressing children without hiding their child deficits.  It uses
one fresh selector and makes no overlapping-input or reused-selector claim.
It is consistent with selector-erasure obstructions for arbitrary two-branch
splices and with the failure of overlapping-product subadditivity.

